\textbf{Phần 1. Trắc nghiệm lựa chọn}
\setcounterpageref{ex}{0}

\begin{ex} %[5X11V65]
Một xạ thủ bắn $80$ viên đạn vào bia, thấy có $68$ viên trúng mục tiêu. Tỉ số phần trăm số lần bắn trượt mục tiêu của xạ thủ đó là
\choice
{$12\%$}
{\True $15\%$}
{$18\%$}
{$85\%$}
\loigiai{
Số lần bắn trượt mục tiêu là:
$$80 - 68 = 12\text{ (lần)}.$$
Tỉ số phần trăm số lần bắn trượt là:
$$12 : 80 \times 100\% = 15\%.$$
\textbf{Đáp án của câu hỏi là B}
}
\end{ex}

\begin{ex} %[5X11V65]
Gieo một con xúc xắc $60$ lần liên tiếp, kết quả ghi lại số lần xuất hiện mặt chẵn ($2, 4, 6$) là $33$ lần. Tỉ số của số lần xuất hiện mặt lẻ so với tổng số lần gieo viết dưới dạng số thập phân là
\choice
{$0{,}55$}
{\True $0{,}45$}
{$0{,}4$}
{$0{,}5$}
\loigiai{
Số lần xuất hiện mặt lẻ là:
$$60 - 33 = 27\text{ (lần)}.$$
Tỉ số số lần xuất hiện mặt lẻ so với tổng số lần gieo là:
$$27 : 60 = 0{,}45.$$
\textbf{Đáp án của câu hỏi là B}
}
\end{ex}

\begin{ex} %[5X11V65]
Trong một hộp có $3$ quả bóng màu: Đỏ, Xanh, Vàng. Bạn Nam lấy ngẫu nhiên một quả bóng, ghi lại màu rồi bỏ lại vào hộp. Sau $50$ lần thực hiện, bóng đỏ xuất hiện $18$ lần, bóng xanh xuất hiện $17$ lần. Tỉ số phần trăm số lần xuất hiện bóng vàng là
\choice
{$25\%$}
{\True $30\%$}
{$35\%$}
{$40\%$}
\loigiai{
Số lần lấy được bóng màu vàng là:
$$50 - (18 + 17) = 15\text{ (lần)}.$$
Tỉ số phần trăm xuất hiện bóng vàng là:
$$15 : 50 \times 100\% = 30\%.$$
\textbf{Đáp án của câu hỏi là B}
}
\end{ex}

\begin{ex} %[5X11V65]
Tung một đồng xu liên tiếp $120$ lần, thấy mặt Sấp xuất hiện với tỉ số số lần lặp lại so với tổng số lần tung là $\dfrac{7}{12}$. Số lần xuất hiện mặt Ngửa là
\choice
{$45\text{ lần}$}
{\True $50\text{ lần}$}
{$55\text{ lần}$}
{$70\text{ lần}$}
\loigiai{
Số lần xuất hiện mặt Sấp là:
$$120 \times \dfrac{7}{12} = 70\text{ (lần)}.$$
Số lần xuất hiện mặt Ngửa là:
$$120 - 70 = 50\text{ (lần)}.$$
\textbf{Đáp án của câu hỏi là B}
}
\end{ex}

\begin{ex} %[5X11V65]
Một cầu thủ thực hiện $40$ quả đá luân lưu, tỉ lệ sút thành công là $85\%$. Số quả sút không thành công của cầu thủ đó là
\choice
{$4\text{ quả}$}
{\True $6\text{ quả}$}
{$8\text{ quả}$}
{$10\text{ quả}$}
\loigiai{
Tỉ số phần trăm số quả sút không thành công là:
$$100\% - 85\% = 15\%.$$
Số quả sút không thành công là:
$$40 \times 15\% = 6\text{ (quả)}.$$
\textbf{Đáp án của câu hỏi là B}
}
\end{ex}

\begin{ex} %[5X11V65]
Quay một vòng quay may mắn chia đều thành $8$ phần bằng nhau gồm $3$ ô trúng thưởng và $5$ ô mất lượt. Bạn An quay $40$ lần và ghi lại kết quả có $18$ lần quay vào ô trúng thưởng. Tỉ số giữa số lần quay trúng thưởng thực nghiệm của An và số lần trúng thưởng theo lí thuyết ($40 \times \dfrac{3}{8}$) là
\choice
{$1$}
{\True $1{,}2$}
{$1{,}25$}
{$1{,}5$}
\loigiai{
Số lần trúng thưởng theo lí thuyết là:
$$40 \times \dfrac{3}{8} = 15\text{ (lần)}.$$
Tỉ số giữa số lần trúng thưởng thực nghiệm và lí thuyết là:
$$18 : 15 = 1{,}2.$$
\textbf{Đáp án của câu hỏi là B}
}
\end{ex}

\begin{ex} %[5X11V65]
Kiểm tra chất lượng $200$ sản phẩm của một dây chuyền sản xuất thấy có $6$ sản phẩm bị lỗi. Tỉ số phần trăm sản phẩm đạt tiêu chuẩn chất lượng của dây chuyền đó là
\choice
{$96\%$}
{\True $97\%$}
{$98\%$}
{$98{,}5\%$}
\loigiai{
Số sản phẩm đạt tiêu chuẩn chất lượng là:
$$200 - 6 = 194\text{ (sản phẩm)}.$$
Tỉ số phần trăm sản phẩm đạt chuẩn là:
$$194 : 200 \times 100\% = 97\%.$$
\textbf{Đáp án của câu hỏi là B}
}
\end{ex}

\begin{ex} %[5X11V65]
Bạn Bình gieo một con xúc xắc $90$ lần và thống kê số lần xuất hiện các mặt từ $1$ đến $6$. Biết mặt $1$ xuất hiện $14$ lần, mặt $2$ xuất hiện $16$ lần, mặt $3$ xuất hiện $15$ lần, mặt $4$ xuất hiện $13$ lần, mặt $5$ xuất hiện $17$ lần. Tỉ số số lần lặp lại của sự kiện “xuất hiện mặt $6$ chấm” so với tổng số lần gieo viết dưới dạng phân số tối giản là
\choice
{$\dfrac{1}{5}$}
{\True $\dfrac{1}{6}$}
{$\dfrac{2}{15}$}
{$\dfrac{7}{45}$}
\loigiai{
Số lần xuất hiện mặt $6$ chấm là:
$$90 - (14 + 16 + 15 + 13 + 17) = 90 - 75 = 15\text{ (lần)}.$$
Tỉ số của số lần xuất hiện mặt $6$ chấm so với tổng số lần gieo là:
$$\dfrac{15}{90} = \dfrac{1}{6}.$$
\textbf{Đáp án của câu hỏi là B}
}
\end{ex}

\begin{ex} %[5X11V65]
Một học sinh làm một bài kiểm tra trắc nghiệm gồm $20$ câu. Bạn đó chọn ngẫu nhiên các câu trả lời và đúng được $6$ câu. Tỉ số phần trăm số câu trả lời sai của học sinh đó là
\choice
{$65\%$}
{\True $70\%$}
{$75\%$}
{$80\%$}
\loigiai{
Số câu trả lời sai là:
$$20 - 6 = 14\text{ (câu)}.$$
Tỉ số phần trăm câu trả lời sai là:
$$14 : 20 \times 100\% = 70\%.$$
\textbf{Đáp án của câu hỏi là B}
}
\end{ex}

\begin{ex} %[5X11V65]
Rút ngẫu nhiên một thẻ từ hộp có $10$ thẻ đánh số từ $1$ đến $10$, ghi lại số rồi trả thẻ vào hộp. Thực hiện $50$ lần rút, sự kiện “rút được thẻ mang số nguyên tố” ($2, 3, 5, 7$) xuất hiện $22$ lần. Tỉ số số lần rút được thẻ không phải số nguyên tố là
\choice
{$\dfrac{11}{25}$}
{\True $\dfrac{14}{25}$}
{$\dfrac{13}{25}$}
{$\dfrac{12}{25}$}
\loigiai{
Số lần rút được thẻ không phải số nguyên tố là:
$$50 - 22 = 28\text{ (lần)}.$$
Tỉ số là:
$$\dfrac{28}{50} = \dfrac{14}{25}.$$
\textbf{Đáp án của câu hỏi là B}
}
\end{ex}

\begin{ex} %[5X11V65]
Người ta theo dõi $120$ chuyến xe buýt đến bến xe và ghi nhận có $102$ chuyến đến đúng giờ. Tỉ số phần trăm số chuyến xe buýt đến muộn giờ là
\choice
{$12\%$}
{\True $15\%$}
{$18\%$}
{$20\%$}
\loigiai{
Số chuyến xe buýt đến muộn giờ là:
$$120 - 102 = 18\text{ (chuyến)}.$$
Tỉ số phần trăm chuyến xe đến muộn là:
$$18 : 120 \times 100\% = 15\%.$$
\textbf{Đáp án của câu hỏi là B}
}
\end{ex}

\begin{ex} %[5X11V65]
Hai bạn Mai và Lan cùng chơi trò tung đồng xu. Mai tung $40$ lần có $22$ lần xuất hiện mặt Ngửa. Lan tung $60$ lần có $28$ lần xuất hiện mặt Ngửa. Khẳng định nào sau đây là đúng?
\choice
{Lan có tỉ lệ xuất hiện mặt Ngửa cao hơn Mai}
{\True Mai có tỉ lệ xuất hiện mặt Ngửa cao hơn Lan}
{Tỉ lệ xuất hiện mặt Ngửa của hai bạn bằng nhau}
{Tổng số lần xuất hiện mặt Sấp của Mai nhiều hơn của Lan}
\loigiai{
Tỉ lệ mặt Ngửa của Mai là:
$$\dfrac{22}{40} = 55\%.$$
Tỉ lệ mặt Ngửa của Lan là:
$$\dfrac{28}{60} \approx 46{,}67\%.$$
Vì $55\% > 46{,}67\%$ nên Mai có tỉ lệ xuất hiện mặt Ngửa cao hơn Lan.\\
\textbf{Đáp án của câu hỏi là B}
}
\end{ex}

\textbf{Phần 2. Trắc nghiệm đúng sai}
\setcounterpageref{ex}{0}

\begin{ex} %[5X11V65]
Một hộp có các thẻ ghi số từ $1$ đến $5$. Bạn Nam rút ngẫu nhiên $1$ thẻ, ghi lại số rồi trả lại hộp. Sau $100$ lần thực hiện thu được bảng kết quả:
- Số 1 xuất hiện $18$ lần;
- Số 2 xuất hiện $24$ lần;
- Số 3 xuất hiện $16$ lần;
- Số 4 xuất hiện $22$ lần;
- Số 5 xuất hiện $20$ lần.
Xét tính đúng/sai của các nhận định sau:
\choiceTF
{\True Số 2 có số lần xuất hiện nhiều nhất trong $100$ lần rút}
{\True Tỉ số phần trăm số lần xuất hiện số lẻ ($1, 3, 5$) là $54\%$}
{Tỉ số phần trăm số lần xuất hiện số 4 là $20\%$}
{\True Số lần xuất hiện số chẵn ít hơn số lần xuất hiện số lẻ là $8$ lần}
\loigiai{
\begin{itemchoice}
\item Đúng. Số 2 xuất hiện 24 lần (nhiều nhất).
\item Đúng. Tổng số lần xuất hiện số lẻ: $18 + 16 + 20 = 54$ lần $\Rightarrow 54\%$.
\item Sai. Số 4 xuất hiện 22 lần $\Rightarrow 22\%$ (không phải 20\%).
\item Đúng. Tổng số lần số chẵn là $24 + 22 = 46$ lần. Chênh lệch: $54 - 46 = 8$ lần.
\end{itemchoice}
}
\end{ex}

\begin{ex} %[5X11V65]
Một xạ thủ bắn súng thể thao bắn $50$ viên đạn vào bia. Kết quả thu được: $35$ viên đạt điểm $10$, $10$ viên đạt điểm $9$, $4$ viên đạt điểm $8$ và $1$ viên đạt điểm $7$. Xét tính đúng/sai của các nhận định sau:
\choiceTF
{\True Tỉ số số lần bắn đạt điểm 10 so với tổng số lần bắn là $\dfrac{7}{10}$}
{\True Tỉ số phần trăm số lần bắn đạt điểm 10 là $70\%$}
{Số lần bắn đạt điểm từ 9 trở lên chiếm $80\%$ tổng số lần bắn}
{\True Số lần bắn đạt điểm dưới 9 chiếm $10\%$ tổng số lần bắn}
\loigiai{
\begin{itemchoice}
\item Đúng. $\dfrac{35}{50} = \dfrac{7}{10}$.
\item Đúng. $\dfrac{35}{50} = 70\%$.
\item Sai. Số lần đạt điểm từ 9 trở lên: $35 + 10 = 45$ lần $\Rightarrow 45 : 50 = 90\%$ (chứ không phải 80\%).
\item Đúng. Số lần bắn dưới 9 điểm: $4 + 1 = 5$ lần $\Rightarrow 5 : 50 = 10\%$.
\end{itemchoice}
}
\end{ex}

\textbf{Phần 3. Câu trả lời ngắn}
\setcounterpageref{ex}{0}

\begin{ex} %[5X11V65]
Tung một đồng xu $80$ lần, thấy có $44$ lần xuất hiện mặt Sấp. Hỏi tỉ số phần trăm số lần xuất hiện mặt Ngửa là bao nhiêu phần trăm?
\shortans{45}
\loigiai{
Số lần xuất hiện mặt Ngửa là:
$$80 - 44 = 36\text{ (lần)}.$$
Tỉ số phần trăm mặt Ngửa là:
$$36 : 80 \times 100\% = 45\%.$$
}
\end{ex}

\begin{ex} %[5X11V65]
Gieo một con xúc xắc $75$ lần, sự kiện “xuất hiện mặt có số chấm chia hết cho $3$” ($3$ hoặc $6$ chấm) xuất hiện $24$ lần. Hỏi tỉ số phần trăm số lần xuất hiện mặt không chia hết cho $3$ là bao nhiêu phần trăm?
\shortans{68}
\loigiai{
Số lần xuất hiện mặt không chia hết cho $3$ là:
$$75 - 24 = 51\text{ (lần)}.$$
Tỉ số phần trăm là:
$$51 : 75 \times 100\% = 68\%.$$
}
\end{ex}

\begin{ex} %[5X11V65]
Một cầu thủ thực hiện sút phạt đền $50$ lần và có tỉ lệ thành công là $88\%$. Hỏi cầu thủ đó đã sút hỏng bao nhiêu quả phạt đền?
\shortans{6}
\loigiai{
Tỉ lệ sút hỏng là:
$$100\% - 88\% = 12\%.$$
Số quả phạt đền sút hỏng là:
$$50 \times 12\% = 6\text{ (quả)}.$$
}
\end{ex}

\begin{ex} %[5X11V65]
Trong một đợt kiểm tra chất lượng sản phẩm bóng đèn, người ta thử nghiệm ngẫu nhiên $500$ chiếc và thấy có $485$ chiếc hoạt động tốt. Hỏi có bao nhiêu bóng đèn bị hỏng?
\shortans{15}
\loigiai{
Số bóng đèn bị hỏng là:
$$500 - 485 = 15\text{ (bóng)}.$$
}
\end{ex}

\begin{ex} %[5X11V65]
Quay một đĩa số tròn $60$ lần, mũi tên dừng ở số may mắn chiếm tỉ lệ $35\%$. Hỏi mũi tên đã dừng ở các số khác bao nhiêu lần?
\shortans{39}
\loigiai{
Tỉ lệ mũi tên dừng ở các số khác là:
$$100\% - 35\% = 65\%.$$
Số lần mũi tên dừng ở các số khác là:
$$60 \times 65\% = 39\text{ (lần)}.$$
}
\end{ex}

\begin{ex} %[5X11V65]
Một nhóm học sinh gieo xúc xắc $120$ lần và ghi nhận mặt $1$ chấm xuất hiện $18$ lần. Hỏi số lần xuất hiện các mặt còn lại (từ $2$ đến $6$) là bao nhiêu lần?
\shortans{102}
\loigiai{
Số lần xuất hiện các mặt còn lại là:
$$120 - 18 = 102\text{ (lần)}.$$
}
\end{ex}
